\documentclass[12pt, a4paper]{article}
\usepackage[T2A]{fontenc}
\usepackage[utf8]{inputenc}
\usepackage[russian]{babel}
\usepackage{amsmath, amssymb} % Для математических символов
\usepackage{enumitem} % Для гибкой настройки списков
\usepackage{geometry} % Для настройки полей
\geometry{top=2cm, bottom=2cm, left=2.5cm, right=2cm}

\begin{document}

\section*{Борис Мозесович Давидович, Петр Евгеньевич Пушкарь, Юрий Витальевич Чеканов <<Математический анализ в 57-й школе. Четырехгодичный курс>>}

\subsection*{1. Теория множеств}

\subsubsection*{Задача 1}
\textbf{Доказать, что:}
\begin{enumerate}[label=\alph*)]
    \item \(A \subset A\);
    \item если \(A \subset B\) и \(B \subset C\), то \(A \subset C\);
    \item \(A \subset B\) и \(B \subset A\) тогда и только тогда, когда \(A = B\).
\end{enumerate}

\textbf{Доказательство:}
\begin{enumerate}[label=\alph*)]
    \item Возьмём произвольный элемент \(x \in A\). Очевидно, что \(x \in A\). Значит, по определению подмножества, \(A \subset A\).
    \item Пусть \(A \subset B\) и \(B \subset C\). Докажем, что \(A \subset C\). Возьмём произвольный элемент \(x \in A\). Так как \(A \subset B\), то \(x \in B\). Так как \(B \subset C\), то \(x \in C\). Мы показали, что любой элемент из \(A\) принадлежит \(C\), значит \(A \subset C\).
    \item 
    \textit{Необходимость.} Пусть \(A \subset B\) и \(B \subset A\). По определению равенства множеств (Определение 1) это означает, что \(A = B\). \\
    \textit{Достаточность.} Пусть \(A = B\). Тогда каждый элемент множества \(A\) является элементом множества \(B\), то есть \(A \subset B\). Аналогично, каждый элемент множества \(B\) является элементом \(A\), то есть \(B \subset A\).
\end{enumerate}

\subsubsection*{Задача 2}
\textbf{Доказать, что:}
\begin{enumerate}[label=\alph*)]
    \item \(\varnothing \subset A\);
    \item если \(A \subset \varnothing\), то \(A = \varnothing\).
\end{enumerate}

\textbf{Доказательство:}
\begin{enumerate}[label=\alph*)]
    \item По определению \(X \subset Y\), если любой элемент из \(X\) принадлежит \(Y\). Так как в пустом множестве \(\varnothing\) нет элементов, то условие «если \(x \in \varnothing\), то \(x \in A\)» выполняется автоматически (из ложного утверждения следует любое). Значит, \(\varnothing \subset A\).
    \item Пусть \(A \subset \varnothing\). Это значит, что любой элемент из \(A\) принадлежит \(\varnothing\). Но в \(\varnothing\) нет элементов, следовательно, в \(A\) также не может быть элементов. Значит, \(A\) — пустое множество, то есть \(A = \varnothing\).
\end{enumerate}

\subsubsection*{Задача 3}
\textbf{Сколько элементов и подмножеств у каждого из следующих множеств:}
а) \(\{0\}\); б) \(\varnothing\); в) \(\{1, 2\}\); г) \(\{1, 2, 3\}\); д) \(\{\varnothing\}\); е) \(\{\{1, 2, 3\}\}\); ж) \(\{\{1, 2\}, 3\}\)?

\textbf{Решение:}
Количество подмножеств множества из \(k\) элементов равно \(2^k\).
\begin{itemize}
    \item а) \(\{0\}\): 1 элемент, \(2^1 = 2\) подмножества (само множество и \(\varnothing\)).
    \item б) \(\varnothing\): 0 элементов, \(2^0 = 1\) подмножество (само \(\varnothing\)).
    \item в) \(\{1, 2\}\): 2 элемента, \(2^2 = 4\) подмножества.
    \item г) \(\{1, 2, 3\}\): 3 элемента, \(2^3 = 8\) подмножеств.
    \item д) \(\{\varnothing\}\): 1 элемент (это само \(\varnothing\)), \(2^1 = 2\) подмножества.
    \item е) \(\{\{1, 2, 3\}\}\): 1 элемент (множество \(\{1, 2, 3\}\)), \(2^1 = 2\) подмножества.
    \item ж) \(\{\{1, 2\}, 3\}\): 2 элемента (множество \(\{1, 2\}\) и число 3), \(2^2 = 4\) подмножества.
\end{itemize}

\subsubsection*{Задача 4}
\textbf{Может ли у множества быть ровно: а) 0; б) 7; в) 16 подмножеств?}

\textbf{Решение:}
Количество подмножеств множества из \(k\) элементов равно \(2^k\).
\begin{itemize}
    \item а) 0: нет, так как \(2^k \ge 1\) (у любого множества есть хотя бы само себя).
    \item б) 7: нет, так как 7 не является степенью двойки.
    \item в) 16: да, так как \(16 = 2^4\) (например, множество из 4 элементов).
\end{itemize}

\subsubsection*{Задача 5}
\textbf{Может ли у множества быть ровно на 2000 подмножеств больше, чем у другого?}

\textbf{Решение:}
Пусть у первого множества \(k\) элементов, у второго \(n\) элементов. Тогда \(2^k - 2^n = 2000\).
Разделим обе части на \(2^n\):
\[
2^{k-n} - 1 = \frac{2000}{2^n}.
\]
Число \(2000 = 2^4 \cdot 125\) содержит только \(2^4\) в качестве степени двойки, значит \(n \le 4\).
Если \(n \le 4\), то \(2^n\) может быть \(1, 2, 4, 8, 16\).
Переберём \(n\):
\begin{itemize}
    \item \(n=4\): \(2^k - 16 = 2000 \Rightarrow 2^k = 2016\) (не степень двойки).
    \item \(n=3\): \(2^k - 8 = 2000 \Rightarrow 2^k = 2008\) (нет).
    \item \(n=2\): \(2^k - 4 = 2000 \Rightarrow 2^k = 2004\) (нет).
    \item \(n=1\): \(2^k - 2 = 2000 \Rightarrow 2^k = 2002\) (нет).
    \item \(n=0\): \(2^k - 1 = 2000 \Rightarrow 2^k = 2001\) (нет).
\end{itemize}
Ни в одном случае \(k\) не является целым. Следовательно, такого быть не может.
\textbf{Ответ:} нет.

\subsubsection*{Задача 6*}
\textbf{Может ли у множества \(A\) быть ровно 2000 подмножеств, не являющихся подмножествами множества \(B\), ни подмножествами множества \(C\)?}

\textbf{Решение:}
Разобьём элементы множества \(A\) на 4 непересекающиеся группы:
\begin{itemize}
    \item \(x\) элементов принадлежат \(B\), но не принадлежат \(C\);
    \item \(y\) элементов принадлежат \(C\), но не принадлежат \(B\);
    \item \(z\) элементов принадлежат и \(B\), и \(C\);
    \item \(w\) элементов не принадлежат ни \(B\), ни \(C\).
\end{itemize}
Общее количество элементов в \(A\) равно \(k = x + y + z + w\).

Подмножество \(X \subset A\) не является подмножеством \(B\), если оно содержит хотя бы один элемент из \(y\) или \(w\). 
Количество таких подмножеств равно \(2^k - 2^{x+z}\) (из всех подмножеств вычитаем те, которые целиком лежат в \(B\)).
Аналогично, количество подмножеств, не являющихся подмножествами \(C\), равно \(2^k - 2^{y+z}\).

Количество подмножеств, которые не являются подмножествами ни \(B\), ни \(C\), по формуле включений-исключений равно:
\[
N = 2^k - (2^k - 2^{x+z}) - (2^k - 2^{y+z}) + (2^k - 2^{x+z} - 2^{y+z} + 2^z).
\]
Упростим выражение:
\[
N = 2^k - 2^{x+z} - 2^{y+z} + 2^z = 2^z(2^{k-z} - 2^x - 2^y + 1).
\]
По условию \(N = 2000\). Так как \(2000 = 2^4 \cdot 125\), то \(2^z\) должно делить 2000, откуда \(z \le 4\).

Пусть \(z=4\). Тогда \(2^4 = 16\), и скобка должна быть равна \(2000/16 = 125\):
\[
2^{k-4} - 2^x - 2^y + 1 = 125 \implies 2^{k-4} - 2^x - 2^y = 124.
\]
Подберём степени двойки. Если \(k-4 = 7\) (то есть \(k=11\)), то \(128 - 2^x - 2^y = 124\), откуда \(2^x + 2^y = 4\). Это выполняется при \(x=1, y=1\).
Тогда \(w = k - (x+y+z) = 11 - (1+1+4) = 5\).

Значит, если в множестве \(A\) 11 элементов, а множества \(B\) и \(C\) состоят из 5 элементов каждое и пересекаются по 4 элементам, то количество подмножеств \(A\), не являющихся подмножествами ни \(B\), ни \(C\), равно:
\[
2^{11} - 2^5 - 2^5 + 2^4 = 2048 - 32 - 32 + 16 = 2000.
\]

\textbf{Ответ:} да, может. Например, \(|A|=11\), \(|B|=|C|=5\), \(|B \cap C|=4\).

\subsubsection*{Задача 7}
\textbf{Пусть \(A = \{1, 3, 7, 137\}\), \(B = \{3, 7, 100\}\), \(C = \{0, 1, 3, 100\}\), \(D = \{0, 7, 100, 333\}\). Найдите множества:}

\textbf{Решение:}
\begin{enumerate}[label=\alph*)]
    \item \(A \cup B = \{1, 3, 7, 100, 137\}\)
    \item \(A \cap B = \{3, 7\}\)
    \item \((A \cap B) \cup D = \{3, 7\} \cup \{0, 7, 100, 333\} = \{0, 3, 7, 100, 333\}\)
    \item \(C \cap (D \cap B) = C \cap (\{0, 7, 100, 333\} \cap \{3, 7, 100\}) = C \cap \{7, 100\} = \{0, 1, 3, 100\} \cap \{7, 100\} = \{100\}\)
    \item \((A \cup B) \cap (C \cup D) = \{1, 3, 7, 100, 137\} \cap \{0, 1, 3, 7, 100, 333\} = \{1, 3, 7, 100\}\)
    \item \((A \cap B) \cup (C \cap D) = \{3, 7\} \cup (\{0, 1, 3, 100\} \cap \{0, 7, 100, 333\}) = \{3, 7\} \cup \{0, 100\} = \{0, 3, 7, 100\}\)
    \item \((D \cup A) \cap (C \cup B) = \{0, 1, 7, 100, 137, 333\} \cap \{0, 1, 3, 7, 100\} = \{0, 1, 7, 100\}\)
    \item \((A \cap (B \cap C)) \cap D = A \cap \{3\} \cap D = \{3\} \cap \{0, 7, 100, 333\} = \varnothing\)
    \item \((A \cup (B \cap C)) \cap D = (A \cup \{3, 100\}) \cap D = \{1, 3, 7, 100, 137\} \cap \{0, 7, 100, 333\} = \{7, 100\}\)
    \item \((C \cap A) \cup ((A \cup (C \cap D)) \cap B) = \{1, 3\} \cup ((\{1, 3, 7, 137\} \cup \{0, 100\}) \cap \{3, 7, 100\}) = \{1, 3\} \cup (\{0, 1, 3, 7, 100, 137\} \cap \{3, 7, 100\}) = \{1, 3\} \cup \{3, 7, 100\} = \{1, 3, 7, 100\}\)
\end{enumerate}

\subsubsection*{Задача 8}
\textbf{Доказать, что для любых множеств \(A, B, C\)}
\begin{enumerate}[label=\alph*)]
    \item \(A \cup A = A\), \(A \cap A = A\);
    \item \(A \cup B = B \cup A\), \(A \cap B = B \cap A\);
    \item \(A \cup (B \cup C) = (A \cup B) \cup C\), \(A \cap (B \cap C) = (A \cap B) \cap C\);
    \item \(A \cap (B \cup C) = (A \cap B) \cup (A \cap C)\), \(A \cup (B \cap C) = (A \cup B) \cap (A \cup C)\).
\end{enumerate}

\textbf{Доказательство:}
\begin{enumerate}[label=\alph*)]
    \item 
    Докажем \(A \cup A = A\). 
    \textit{1) Покажем, что \(A \cup A \subset A\).} Пусть \(x \in A \cup A\). По определению объединения, \(x \in A\) или \(x \in A\). В любом случае \(x \in A\). Следовательно, \(A \cup A \subset A\). \\
    \textit{2) Покажем, что \(A \subset A \cup A\).} Пусть \(x \in A\). Тогда истинно высказывание «\(x \in A\) или \(x \in A\)», что означает \(x \in A \cup A\). Следовательно, \(A \subset A \cup A\). \\
    Из двух включений следует \(A \cup A = A\). 
    Равенство \(A \cap A = A\) доказывается аналогично (слово «или» заменяется на «и»).

    \item 
    Докажем \(A \cup B = B \cup A\) (коммутативность объединения).
    \textit{1) Покажем, что \(A \cup B \subset B \cup A\).} Пусть \(x \in A \cup B\). По определению \(x \in A\) или \(x \in B\). 
    Если \(x \in A\), то истинно «\(x \in B\) или \(x \in A\)», значит \(x \in B \cup A\).
    Если \(x \in B\), то аналогично \(x \in B \cup A\).
    В любом случае \(x \in B \cup A\), значит \(A \cup B \subset B \cup A\). \\
    \textit{2) Покажем, что \(B \cup A \subset A \cup B\).} Пусть \(x \in B \cup A\). Рассуждая аналогично, получаем, что \(x \in A \cup B\).
    Значит, \(B \cup A \subset A \cup B\). \\
    Следовательно, \(A \cup B = B \cup A\).
    Равенство \(A \cap B = B \cap A\) доказывается аналогично.

    \item 
    Докажем \(A \cup (B \cup C) = (A \cup B) \cup C\) (ассоциативность).
    \textit{1) \(A \cup (B \cup C) \subset (A \cup B) \cup C\).} Пусть \(x \in A \cup (B \cup C)\). Тогда \(x \in A\) или \(x \in (B \cup C)\).
    \begin{itemize}
        \item Если \(x \in A\), то \(x \in A \cup B\), а значит, \(x \in (A \cup B) \cup C\).
        \item Если \(x \in (B \cup C)\), то \(x \in B\) или \(x \in C\). Если \(x \in B\), то \(x \in A \cup B\), следовательно \(x \in (A \cup B) \cup C\). Если \(x \in C\), то по определению объединения \(x \in (A \cup B) \cup C\).
    \end{itemize}
    В любом случае включение доказано. \\
    \textit{2) \((A \cup B) \cup C \subset A \cup (B \cup C)\).} Пусть \(x \in (A \cup B) \cup C\). Тогда \(x \in (A \cup B)\) или \(x \in C\). Если \(x \in C\), то \(x \in B \cup C\), значит \(x \in A \cup (B \cup C)\). Если \(x \in A \cup B\), то \(x \in A\) или \(x \in B\). Если \(x \in A\), то включение очевидно. Если \(x \in B\), то \(x \in B \cup C\), следовательно \(x \in A \cup (B \cup C)\).
    Включение доказано. \\
    Второе равенство \(A \cap (B \cap C) = (A \cap B) \cap C\) доказывается аналогично.

    \item 
    Докажем дистрибутивность \(A \cap (B \cup C) = (A \cap B) \cup (A \cap C)\).
    \textit{1) \(A \cap (B \cup C) \subset (A \cap B) \cup (A \cap C)\).} Пусть \(x \in A \cap (B \cup C)\). Тогда \(x \in A\) и \(x \in (B \cup C)\). Значит, \(x \in B\) или \(x \in C\).
    Если \(x \in B\), то \(x \in A\) и \(x \in B\), значит \(x \in A \cap B\), откуда \(x \in (A \cap B) \cup (A \cap C)\).
    Если \(x \in C\), то \(x \in A\) и \(x \in C\), значит \(x \in A \cap C\), откуда \(x \in (A \cap B) \cup (A \cap C)\).
    Включение доказано. \\
    \textit{2) \((A \cap B) \cup (A \cap C) \subset A \cap (B \cup C)\).} Пусть \(x \in (A \cap B) \cup (A \cap C)\). Тогда \(x \in A \cap B\) или \(x \in A \cap C\).
    В любом из этих случаев \(x \in A\). 
    Если \(x \in A \cap B\), то \(x \in B\), значит \(x \in B \cup C\). Отсюда \(x \in A \cap (B \cup C)\).
    Если \(x \in A \cap C\), то \(x \in C\), значит \(x \in B \cup C\). Отсюда \(x \in A \cap (B \cup C)\).
    Включение доказано. \\
    Значит, равенство верно. 
    Второе равенство \(A \cup (B \cap C) = (A \cup B) \cap (A \cup C)\) доказывается аналогично.
\end{enumerate}

\end{document}